\section{Hybrid Architectures and Controls}
\label{sec:hybrid}

\subsection{Motivation and Architectural Formulation (H1)}
While Message Passing Neural Networks (MPNNs)~\cite{gilmer2017neural} effectively capture continuous, multi-hop topological features from molecular graphs, their local message aggregation mechanism exhibits an inherent limitation: whole-molecule physicochemical invariants---such as lipophilicity ($\log P$), topological polar surface area (TPSA), molecular weight, and formal charge distributions---can only be approximated through extensive message passing iterations. On small drug discovery benchmarks such as Tox21 ($\sim$6,200 training instances), deep propagation risks oversmoothing and gradient degradation. Supplying low-dimensional, preprocessed physicochemical descriptors provides an informative global prior that does not require re-discovery from graph topology.

To exploit this structural complementarity, we formulate the H1 hybrid architecture (\texttt{mpnn+desc}). The graph encoder processes node features $\mathbf{x}_v \in \mathbb{R}^{33}$ and directed edge features $\mathbf{e}_{vw} \in \mathbb{R}^6$ through an edge-network message passing phase with GRU updates across $T=3$ steps, pooling all atom states via an order-invariant Set2Set readout into a graph embedding $\mathbf{z}_{\text{graph}} \in \mathbb{R}^{128}$ ($2 \times 64$). Concurrently, the 213 preprocessed, leakage-free 2D RDKit descriptors $\mathbf{d} \in \mathbb{R}^{213}$ pass through an auxiliary dense projection layer $\mathbf{z}_{\text{desc}} = \text{Dropout}(\text{ReLU}(\mathbf{W}_d \mathbf{d} + \mathbf{b}_d))$ with $d_{\text{hidden}} = 128$. The representations are fused via concatenation $\mathbf{z} = [\mathbf{z}_{\text{graph}} \,\|\, \mathbf{z}_{\text{desc}}] \in \mathbb{R}^{256}$ before reaching a two-layer multi-task classification head mapping to 12 endpoint logits. Training is optimized using masked, class-weighted binary cross-entropy.

\subsection{Empirical Validation and Control Baselines}
A key vulnerability in hybrid deep learning is attributing performance increases to feature synergy when they in fact stem from added parameter capacity or regularization. We strictly benchmark H1 on the primary Bemis-Murcko scaffold split against rigorous controls:
\begin{enumerate}
    \item \textbf{FP-MLP Control (\texttt{fpmlp+both})}: To determine whether performance gains reflect neural architecture vs.\ graph message passing, a multi-layer perceptron was trained directly on concatenated Morgan fingerprints and RDKit descriptors (2,261 inputs). The FP-MLP achieves a scaffold test PR~AUC of $0.3429 \pm 0.0106$ and ROC~AUC of $0.7269 \pm 0.0051$. H1 (\texttt{mpnn+desc}) attains $0.3553 \pm 0.0057$ PR~AUC and $0.7384 \pm 0.0066$ ROC~AUC, outperforming FP-MLP by $+0.0124$ PR~AUC and winning on 9 of the 12 endpoints ($p = 0.151$).
    \item \textbf{Shuffled Descriptors Negative Control (\texttt{shuffled\_desc})}: To verify that the descriptor branch injects genuine molecule-specific information, descriptor rows were randomly permuted across molecules prior to H1 ingestion. Under identical parameters and feature marginal distributions, performance collapses dramatically to PR~AUC $0.2547 \pm 0.0025$ and ROC~AUC $0.6951 \pm 0.0049$ (a severe penalty of $-0.1006$ PR~AUC vs.\ unablated H1). This proves that the performance of H1 is strictly driven by authentic compound-property alignment.
    \item \textbf{Feature Substitution Ablation}: Substituting descriptors with 2,048-bit Morgan fingerprints (\texttt{mpnn+morgan}) drops PR~AUC to $0.2614 \pm 0.0088$, while combining both (\texttt{mpnn+both}) yields $0.3309 \pm 0.0089$. High-dimensional, sparse local bits dilute the gradient of the graph encoder, confirming that compact, continuous global descriptors are the optimal complement to graph message passing.
\end{enumerate}

\subsection{Late-Fusion Ensemble Hybrid (H2)}
To exploit diverse inductive biases without parameter co-adaptation or risk of data leakage, we implement the H2 ensemble (\texttt{ens\_xgb\_mpnn+desc}) by rank-averaging the predictions of the best Phase~1 tree model (XGBoost on \texttt{morgan+desc}) and the deep learning hybrid H1 (\texttt{mpnn+desc}). Ranks are normalized endpoint-wise to $(0, 1]$ and combined with equal weights ($0.5/0.5$) fixed a priori:
\begin{equation}
\hat{P}_{\text{ens}, j} = \frac{1}{2} \left( \frac{\text{rank}(\hat{P}_{\text{xgb}, j})}{N_{\text{test}}} + \frac{\text{rank}(\hat{P}_{\text{h1}, j})}{N_{\text{test}}} \right)
\end{equation}
On the scaffold split, H2 achieves PR~AUC $0.3945 \pm 0.0011$ and ROC~AUC $0.7643 \pm 0.0017$. H2 surpasses H1 across all 12 individual assay endpoints ($12/0$ win/loss, $p = 4.88 \times 10^{-4}$, mean gain $+0.0392$) while establishing the lowest seed variance among all evaluated models ($\sigma_{\text{PR}} = 0.0011$). On the random split, H2 reaches PR~AUC $0.5626 \pm 0.0034$ and ROC~AUC $0.8604 \pm 0.0036$, establishing an effective synergy between gradient-boosted decision trees and graph representation learning.
